\section*{Hoofdstuk \ref{ch:b3:sequence-stratigraphy} --- Sequentiestratigrafie}

\begin{solution}{exo:b3:sequence-stratigraphy:1}
$A = 0.5 + 0.1 = \qty{0.6}{m/kyr}$; in \qty{20}{kyr} is dat \qty{12}{m} aan
accommodatie. Met een dalende zeespiegel geldt $A = -0.3 + 0.1 = \qty{-0.2}{m/kyr}$: een relatieve daling,
afgezwakt door de subsidentie.
\end{solution}

\begin{solution}{exo:b3:sequence-stratigraphy:2}
(a) \omterm{def:b3:sequence-stratigraphy:tract}{Systeemtract van de dalende zeespiegel} (\omterm{def:b3:sequence-stratigraphy:regression}{geforceerde regressie}). (b) \omterm{def:b3:sequence-stratigraphy:tract}{Transgressief systeemtract}. (c)
\omterm{def:b3:sequence-stratigraphy:tract}{Hoogstandsysteemtract}. (d) \omterm{def:b3:sequence-stratigraphy:tract}{Laagstandsysteemtract} (\omterm{def:b3:sequence-stratigraphy:regression}{normale regressie} na de daling).
\end{solution}

\begin{solution}{exo:b3:sequence-stratigraphy:3}
Elke \omterm{def:b3:sequence-stratigraphy:parasequence}{parasequentie} eindigt ondieper dan die eronder: een progradationele stapeling,
waarbij de kustlijn zeewaarts verschuift, zoals bij een \omterm{def:b3:sequence-stratigraphy:tract}{hoogstandsysteemtract} of een \omterm{def:b3:sequence-stratigraphy:tract}{laagstandsysteemtract}. De
omgekeerde volgorde, waarbij elke \omterm{def:b3:sequence-stratigraphy:parasequence}{parasequentie} fijner en dieper eindigt, is retrogradationeel: een
\omterm{def:b3:sequence-stratigraphy:tract}{transgressief systeemtract}.
\end{solution}

\begin{solution}{exo:b3:sequence-stratigraphy:4}
Vijf couplets per bundel is de verhouding van de korte excentriciteit (ongeveer
\qty{100}{kyr}) tot de precessie (ongeveer \qty{20}{kyr}): elk couplet is een
precessiecyclus, elke bundel een excentriciteitscyclus. De snelheid bedraagt
$1.6/100 = \qty{1.6}{cm/kyr}$, en een couplet van \qty{0.32}{m} duurt ongeveer
\qty{20}{kyr}, wat klopt.
\end{solution}

\begin{solution}{exo:b3:sequence-stratigraphy:5}
$\mathrm dx_s/\mathrm dt = (4000 - 0.15\times20\,000)/50 = \qty{20}{m/kyr}$,
zeewaarts. De kustlijn staat stil wanneer $AW = Q$: $A = 4000/20\,000 =
\qty{0.2}{m/kyr}$.
\end{solution}

\begin{solution}{exo:b3:sequence-stratigraphy:6}
$\mathrm dx_s/\mathrm dt = 4000/50 + 0.2/0.01 = 80 + 20 = \qty{100}{m/kyr}$;
de daling alleen verklaart \qty{20}{m/kyr}, een vijfde.
\end{solution}

\begin{solution}{exo:b3:sequence-stratigraphy:7}
Hij trekt zich terug wanneer $Q < AW = 0.4\times25\,000 = \qty{10000}{m^3/kyr}$ per
meter. Met $Q = 6000$: $(6000 - 10\,000)/10 = \qty{-400}{m/kyr}$, een terugtrekking
van \qty{0.4}{km} per duizend jaar: een lage onderwateroever geeft een snelle
transgressie.
\end{solution}

\begin{solution}{exo:b3:sequence-stratigraphy:8}
$1/94.9 - 1/123.9 = \num{0.00247}\text{ cycli}/\unit{kyr}$: een reeks langer dan
$1/0.00247 = \qty{406}{kyr}$ scheidt ze, en \qty{600}{kyr} (resolutie
0.00167) is voldoende. $1/94.9 - 1/98.9 = 0.00042$: dit vereist ongeveer
\qty{2400}{kyr}, vier keer zo lang als de reeks.
\end{solution}

\begin{solution}{exo:b3:sequence-stratigraphy:9}
$16.4/(20\times40.98) = \qty{0.0200}{m/kyr}$, \qty{2.0}{cm/kyr}; de aslagen
liggen $3.1/0.020 = \qty{155}{kyr}$ uit elkaar.
\end{solution}

\begin{solution}{exo:b3:sequence-stratigraphy:10}
(a) $\tan\alpha = -0.4/200 = -0.002$: dalend onder \ang{0.11}. (b)
$0.2/50 = 0.004$: stijgend onder \ang{0.23}. Een seismische doorsnede beslaat tientallen
kilometers bij een verticale overhoging van tientallen tot honderden, dus de hoek is zichtbaar;
een klif van enkele honderden meters lang toont een stijging van een meter of minder, die verloren gaat in
de variatie van de lagen.
\end{solution}

\begin{solution}{exo:b3:sequence-stratigraphy:11}
$1\,296\,000/54.725282 = \qty{23.68}{kyr}$ ($g_5$); $/57.924383 =
\qty{22.37}{kyr}$ ($g_2$); $/68.377912 = \qty{18.95}{kyr}$ ($g_4$);
$/67.834313 = \qty{19.11}{kyr}$ ($g_3$). $g_2 - g_5 = 3.199101''$ per jaar,
\qty{405.1}{kyr}; $g_4 - g_5 = 13.65263''$, \qty{94.9}{kyr}.
\end{solution}

\begin{solution}{exo:b3:sequence-stratigraphy:12}
Beide leveren $0.3\times100 = \qty{30}{m}$ aan accommodatie tegen dezelfde snelheid; met
dezelfde toevoer is de geometrie gelijk, omdat sediment alleen reageert op het
relatieve zeeniveau. Een eustatische verandering is overal gelijk: dezelfde
reeks, van dezelfde ouderdom, op andere randen wijst op de eerste geschiedenis. Een
subsidentiegeschiedenis opgebouwd uit de thermische en belastingsgeschiedenis van het bekken
(backstripping), of een onafhankelijk eustatisch archief zoals zuurstofisotopen,
scheidt ze eveneens.
\end{solution}

\begin{solution}{pb:b3:sequence-stratigraphy:1}
\textbf{1.} $\dot E = 30\times2\pi/400\times\sin(2\pi t/400)$: maximaal
\qty{0.471}{m/kyr}, stijgend bij \qty{100}{kyr} en \qty{500}{kyr}, dalend bij
\qty{300}{kyr}.

\textbf{2.} \qty{0.020}{m/kyr} bij het scharnier, 0.037 op \qty{20}{km}, 0.062 op
\qty{50}{km}, \qty{0.120}{m/kyr} op \qty{120}{km}.

\textbf{3.} $A < 0$ wanneer $\sin(2\pi t/400) < -0.0346/0.471 = -0.0734$: $t =
200 + (400/2\pi)\arcsin0.0734 = \qty{204.7}{kyr}$.

\textbf{4.} Met $\sigma = 0.0617$: $t = 400 - 63.66\arcsin0.131 =
\qty{391.6}{kyr}$.

\textbf{5.} Ongeveer \qty{187}{kyr} tegenover 200: de eustatische daling moet eerst
de subsidentie inhalen en houdt op haar te overtreffen voordat de laagstand is bereikt.

\textbf{6.} $0.045\times400 = \qty{18}{m}$, tegenover een eustatisch bereik van
\qty{60}{m}: een derde.

\textbf{7.} Landwaarts van 50 tot ongeveer \qty{120}{kyr}, zeewaarts tot ongeveer
\qty{440}{kyr}, landwaarts tot ongeveer \qty{560}{kyr}, weer zeewaarts tot
\qty{600}{kyr} (de keerpunten van het model liggen op 118.5, 439.5 en
\qty{559}{kyr}; een tabel om de \qty{50}{kyr} brengt ze alleen in een interval).

\textbf{8.} $(6000 - 0.30\times15\,000)/40 = \qty{37.5}{m/kyr}$, zeewaarts.

\textbf{9.} $6000/15\,000 = \qty{0.40}{m/kyr}$.

\textbf{10.} $6000/35.7 + 0.424/0.012 = 168.1 + 35.3 = \qty{203}{m/kyr}$,
binnen \qty{0.5}{\percent} van het model.

\textbf{11.} De aanvoer: \qty{83}{\percent} van de snelheid.

\textbf{12.} $(-42.7 - 16.3)/(50.7 - 17.4) = \qty{-1.77}{m/km}$, \ang{0.10}.
Sinds zijn tijd is het eerste punt gesubsideerd, $0.0345\times395.5 =
\qty{13.6}{m}$, het tweede $0.0623\times208.5 = \qty{13.0}{m}$: vrijwel
hetzelfde, dus in de huidige hoogteligging verschilt de daling, \qty{59.0}{m},
van de eustatische daling slechts met \qty{0.6}{m}. De relatieve daling aan de
kustlijn op dat moment, eustasie min de subsidentie tijdens de daling, was kleiner.

\textbf{13.} 150 en 200: hoogstand; 250, 300, 350: dalende zeespiegel; 400:
laagstand; 450, 500, 550: transgressief (de kustlijn trekt zich net na
440 terug); 600: hoogstand.

\textbf{14.} 204.7 en \qty{391.6}{kyr} tegenover 204.5 en 391.5: hetzelfde,
binnen de stap van het model van \qty{0.5}{kyr}.

\textbf{15.} Op het plat: het oppervlak dat tijdens de daling bloot kwam te liggen, de bovenkant van
de afzettingen van de hoogstand vanaf het scharnier tot \qty{17}{km}, doorgetrokken over de
naar beneden verspringende voorlaagjes van de dalende zeespiegel; in het bekken de correlatieve
conformiteit aan de basis of de top van de \omterm{def:b3:sequence-stratigraphy:tract}{systeemtract van de dalende zeespiegel}, bij afspraak. Een
echt plat zou door rivieren tot \omterm{def:b3:sequence-stratigraphy:condensed}{ingesneden dalen} worden gesneden, verweerd tot bodems,
en deels geërodeerd zijn.

\textbf{16.} $(18.5 - 52.8)/(559 - 439.5) = \qty{-0.29}{km/kyr}$; het
\omterm{def:b3:sequence-stratigraphy:surfaces}{maximaal overstromingsvlak} vormt zich aan het eind ervan.

\textbf{17.} De kustlijn schuift snel over een bijna vlakke kustvlakte: de
aanvoer blijft er dicht bij gevangen, verspreid over een breed gebied, en elke
plek krijgt weinig. Bekkenwaarts van de teen, voorbij ongeveer \qty{55}{km}, is
de zeebodem uitgehongerd: daar vormt zich een \omterm{def:b3:sequence-stratigraphy:condensed}{gecondenseerde sectie}.

\textbf{18.} Resolutie $1/30 = \num{0.033}\text{ cycli}/\unit{m}$; Nyquist $1/(2\times
0.01) = \num{50}\text{ cycli}/\unit{m}$.

\textbf{19.} $7.5/2.0 = 3.75$, 1 en $0.81/2.0 = 0.405$; de baanverhoudingen
zijn $405/100 = 4.05$ en $41/100 = 0.41$.

\textbf{20.} \qty{2.0}{m} voor \qty{100}{kyr}: \qty{2.0}{cm/kyr}; $0.81/41 =
\qty{1.98}{cm/kyr}$; $7.5/405 = \qty{1.85}{cm/kyr}$. De golflengten die het
\omterm{def:b3:sequence-stratigraphy:periodogram}{periodogram} nabij \qty{8}{m} kan tonen zijn 10, 7.5 en \qty{6}{m}: de werkelijke
\qty{8.1}{m} valt tussen twee ervan.

\textbf{21.} $30/0.020 = \qty{1500}{kyr}$.

\textbf{22.} $23.7\times0.020 = \qty{0.474}{m}$. Bemonsterd elke
\qty{25}{cm}, bedraagt de Nyquistfrequentie 2 cycli per meter, een golflengte van
\qty{0.5}{m}: de precessie, bij 2.11 cycli per meter, ligt erboven en klapt terug
naar $4 - 2.11 = 1.89$, een valse golflengte van \qty{0.53}{m}.

\textbf{23.} Een enkele piek past bij een of andere periode bij een of andere snelheid, ongeacht
zijn oorsprong; drie pieken waarvan de verhoudingen overeenkomen met de baanverhoudingen, zijn onafhankelijk van
de snelheid, die dan bij elk hetzelfde moet uitkomen.

\textbf{24.} De kustlijn verschuift van \qty{17.37}{km} aan het begin van de
daling tot \qty{52.84}{km} aan het eind van de laagstandregressie: een maximale
bekkenwaartse verschuiving van \textbf{\qty[text-series-to-math=true]{35.5}{km}}.
\end{solution}
